\section{Projection Admissibility}

This application demonstrates the minimal conditions under which
a formal relational statement becomes physically admissible.

\subsection{Setup}

Let
\[
\mathcal{U}
\]
denote the relational reference space.

Let
\[
\Pi_C : \mathcal{U}_C \rightarrow m_2(C)
\]
be a projection into an observable calculation domain under an explicit
relational context $C$, including its TRD coupling.

Consider a relational expression
\[
S(u)
\quad\text{with}\quad
u \in \mathcal{U}.
\]

\subsection{Admissibility Condition}

A statement is projection-compatible only if its value depends exclusively
on the projected state within the fixed context $C$.

Formally, projection compatibility requires:
\[
\forall u_1,u_2 \in \mathcal{U} :
\Pi_C(u_1) = \Pi_C(u_2)
\;\Rightarrow\;
S(u_1) = S(u_2).
\]

Physical admissibility additionally requires
\[
u_1,u_2 \in m_2(C),
\]
so that the projected relational origin is uniquely reconstructible
under the same context and projection.

\subsection{Interpretation}

The compatibility condition ensures that the statement does not distinguish
between relational states that are observationally identical.

It is necessary but not sufficient for physical admissibility:
CRA also requires a single explicit validity domain and context,
while $m_2(C)$ requires unique back-reconstruction.

\subsection{Boundary Case}

If there exist
\[
u_1 \neq u_2
\quad\text{such that}\quad
\Pi_C(u_1) = \Pi_C(u_2)
\quad\text{and}\quad
S(u_1) \neq S(u_2),
\]
then the statement crosses a projection boundary.

The computation remains formally valid,
but loses physical interpretability.

\subsection{Result}

Projection compatibility and physical admissibility are distinct
structural criteria.

It does not restrict formal mathematics.
It restricts physical interpretation.

Different explicit contexts may yield different statement values without
contradiction. Ambiguous or mixed contexts are not physically admissible.
